%%%%%%%%%%  scan idx 204  |  folio 191  %%%%%%%%%%
on se
ramène aussitôt au cas où $E\in\Ob\Qcoh(S)$. Le composé
\[
 E\longrightarrow \mathbb RQE\longrightarrow E
\]
des flèches d'adjonction étant l'identité, l'assertion résulte donc de (3.5.1). Ceci achève la démonstration de (3.5).

\textbf{Remarque 3.6.} La conclusion de (3.5) serait fausse si l'on omettait l'hypothèse « séparé ou noethérien », comme le montre un joli contre-exemple de J.L. Verdier (voir Appendice I au présent exposé).

\textbf{Proposition 3.7.} Sous l'hypothèse de (3.5), supposons en outre qu'il existe un entier $N$ tel que l'on ait, pour tout $F\in\Ob\Mod(S)$ et tout ouvert quasi-compact $U$ de $S$,
\[
 H^i(U,F)=0\quad\text{pour }i>N\qquad\text{\footnotemark}\ .
\]
\footnotetext{Condition réalisée, d'après (Tohoku 3.6.5), si l'espace sous-jacent à $S$ est noethérien et de dimension de Krull $\leq N$.}
Alors :
\begin{enumerate}[label=\alph*)]
\item Le cohéreur est de dimension cohomologique finie.
\item Le foncteur
\[
 \varphi:\mathrm D(\Qcoh(S))\longrightarrow \mathrm D(S)
\]
est pleinement fidèle et a pour image essentielle la sous-catégorie pleine $\mathrm D(S)_{\qcoh}$ de $\mathrm D(S)$ formée des complexes à cohomologie quasi-cohérente.
\end{enumerate}

\emph{Preuve de} a). Si $S$ est affine, la formule $Q=\widetilde{\Gamma}_S$ (cf. preuve de (3.2)) implique que $Q$ est de dimension cohomologique $\leq N$. Passons au cas général. Montrons qu'il existe un entier $a$ tel que l'on ait $R^qQ(E)=0$ pour tout $E\in\Ob\Mod(S)$ et tout $q>a$. L'hypothèse implique que tout $E\in\Ob\Mod(S)$ admet une résolution finie droite $E^\bullet$ de longueur $N$ telle que les $E^i$ soient acycliques, i.e. vérifient $H^q(U,E^i)=0$ pour tout ouvert $U$ de $S$ et tout $q>0$. Par suite, il suffit de prouver qu'il existe un entier $b$ tel que $R^qQ(E)=0$ pour tout $\mathcal O_S$-Module acyclique $E$ et tout $q>b$. (On pourra alors prendre $a=b+N$). Soit $(U_i)$ un recouvrement fini de $S$ par $N'$ ouverts
